\chapter{Probability: uncertainty, evidence, and comparison}
\label{learn:probability}

\begin{learningcard}{Prerequisites and question}
Bring fractions, sums, functions, and the logic of Chapter~\ref{learn:proof}.
Calculus becomes useful for continuous distributions, but our main examples
have finite sample spaces. A positive diagnostic result sounds decisive.
How should its meaning depend on the condition's prevalence and the test's
errors? The same reasoning helps interpret a simulation's apparent success.
\end{learningcard}

\section{Specify the experiment before counting}
A sample space $\Omega$ is the set of possible outcomes. An event is a
subset of outcomes whose occurrence we can ask about. On a finite space,
assign probabilities $p(\omega)\geq0$ with
$\sum_{\omega\in\Omega}p(\omega)=1$; the probability of an event is the
sum over its outcomes. Equal probabilities require a modeling assumption,
such as an ideal fair die. Counting favorable outcomes divided by all outcomes
works only under that assumption.

A random variable $X$ is a real-valued function of the outcome. Its expectation
and variance are
\[
 \mathbb E[X]=\sum_{\omega}p(\omega)X(\omega),\qquad
 \operatorname{Var}(X)=\mathbb E[(X-\mathbb E[X])^2].
\]
Expectation is a probability-weighted mean, and variance measures squared
spread about that mean. Expanding the square gives
$\operatorname{Var}(X)=\mathbb E[X^2]-(\mathbb E[X])^2$.
Standard deviation is the nonnegative square root of variance and has the
same units as $X$.

\section{Worked example: a Bernoulli observation}
Let $X$ equal one when a trial succeeds and zero otherwise, with success
probability $p$. Then $\mathbb E[X]=p$, and $X^2=X$ gives
$\operatorname{Var}(X)=p-p^2=p(1-p)$. For independent identically
distributed trials $X_1,\ldots,X_n$, their sample mean
$\overline X=n^{-1}\sum_jX_j$ has expectation $p$. Independence makes
cross covariances zero, so
\[
 \operatorname{Var}(\overline X)=\frac{p(1-p)}n.
\]
The standard deviation of the mean therefore scales as $1/\sqrt n$.
Four times as many independent observations halve it. Repeating the same
stored observation four times creates dependence and supplies no such gain.

More generally, covariance is
$\operatorname{Cov}(X,Y)=\mathbb E[(X-\mathbb E[X])(Y-\mathbb E[Y])]$.
The variance of a sum includes twice the pairwise covariances. Correlation
normalizes covariance by the two positive standard deviations. Independence
implies zero covariance when the moments exist; the converse can fail.
For a uniformly distributed $X$ on $\{-1,0,1\}$, set $Y=X^2$.
Then $\mathbb E[X]=\mathbb E[X^3]=0$, so covariance vanishes, although
$Y$ is determined by $X$.

\section{Worked example: condition on the result you observed}
For events $A,B$ with $P(B)>0$, conditional probability is
$P(A\mid B)=P(A\cap B)/P(B)$. Applying this definition in both orders
yields Bayes's rule:
\[
 P(A\mid B)=\frac{P(B\mid A)P(A)}{P(B)}.
\]
Let $D$ mean a condition is present, and $+$ a positive test. Assume a
population prevalence $P(D)=0.01$, sensitivity $P(+\mid D)=0.99$,
and false-positive probability $P(+\mid D^c)=0.05$. These are hypothetical
teaching parameters, rather than clinical advice about a real test.
Partitioning by condition status gives
\[
 P(+)=0.99(0.01)+0.05(0.99)=0.0594,
 \qquad P(D\mid+)=\frac{0.0099}{0.0594}=\frac16.
\]
A hypothetical population of 10,000 makes the denominator visible: among
100 people with the condition we expect 99 positive results; among 9,900
others we expect 495. Only 99 of the expected 594 positives come from the
condition group. Expected counts may be nonintegers in other examples;
they describe averages under the model rather than guaranteed outcomes.

\begin{figure}[htbp]
  \centering
  \begin{tikzpicture}[x=1cm,y=1cm,>=Stealth,font=\small,line width=0.8pt]
  \node[align=center] (population) at (0,0) {Population};
  \node (condition) at (3.5,1.8) {$D$};
  \node (other) at (3.5,-1.8) {$D^c$};
  \node[anchor=west] (truepositive) at (7,2.7) {$+$: $P(D\cap+)=99/10000$};
  \node[anchor=west] (falsenegative) at (7,0.9) {$-$: $P(D\cap-)=1/10000$};
  \node[anchor=west] (falsepositive) at (7,-0.9) {$+$: $P(D^c\cap+)=495/10000$};
  \node[anchor=west] (truenegative) at (7,-2.7) {$-$: $P(D^c\cap-)=9405/10000$};
  \draw[->] (population)--(condition) node[midway,above,sloped] {$P(D)=1/100$};
  \draw[->] (population)--(other) node[midway,below,sloped] {$P(D^c)=99/100$};
  \draw[->,line width=1.2pt] (condition)--(truepositive.west)
    node[midway,above,sloped] {$P(+\mid D)=99/100$};
  \draw[->,dashed] (condition)--(falsenegative.west)
    node[midway,below,sloped] {$P(-\mid D)=1/100$};
  \draw[->,line width=1.2pt] (other)--(falsepositive.west)
    node[midway,above,sloped] {$P(+\mid D^c)=5/100$};
  \draw[->,dashed] (other)--(truenegative.west)
    node[midway,below,sloped] {$P(-\mid D^c)=95/100$};
  \node[align=center] at (5.8,-3.7)
    {Keep the two positive leaves, then normalize:\\[3pt]
     $\displaystyle P(D\mid+)=\frac{99/10000}{(99+495)/10000}=\frac16$.};
\end{tikzpicture}

  \caption{Read each path from left to right and multiply its conditional branch probabilities to obtain the joint probability at the leaf. Restrict attention to the two solid positive-result branches, add their leaf probabilities, and normalize the condition-positive leaf by that total.}
  \label{fig:learn:bayes_tree}
\end{figure}

\begin{learningcard}{Historical situation: a paper communicated by another reader}
Thomas Bayes's essay appeared in 1763 after his death, communicated by
Richard Price to the Royal Society. The publication's introductory letter
makes Price's role in presenting the work part of the surviving evidence.
The essay investigates inverse-probability reasoning; our diagnostic example
is a modern application rather than a historical episode from the essay.
\footnote{Thomas Bayes, ``An Essay towards solving a Problem in the
Doctrine of Chances,'' communicated by Richard Price,
\emph{Philosophical Transactions} 53 (1763), pp.~370--418, introductory
letter and propositions:
\url{https://doi.org/10.1098/rstl.1763.0053}.}
A named theorem can conceal the editorial and institutional work through
which a mathematical argument reaches its readers.
\end{learningcard}

\section{An interval is a procedure with assumptions}
Suppose observations are independent normal random variables with unknown
mean $\mu$ and known positive standard deviation $\sigma$. Their mean is
normal with standard deviation $\sigma/\sqrt n$. If $z_{0.975}$ denotes
the 97.5th percentile of the standard normal distribution, approximately
$1.96$, then the interval
\[
 \left[\overline X-z_{0.975}\frac{\sigma}{\sqrt n},\,
       \overline X+z_{0.975}\frac{\sigma}{\sqrt n}\right]
\]
has 95 percent coverage: over repeated samples under these assumptions,
95 percent of intervals contain the fixed $\mu$. A particular observed
interval either contains that fixed value or misses it. A Bayesian credible
interval instead assigns posterior probability to parameter values under a
specified prior and likelihood.

Unknown variance, nonnormal small samples, dependence, selection, and repeated
comparisons change the calculation or its interpretation. For independent
normal observations with unknown variance, a Student $t$ interval uses the
sample standard deviation and $n-1$ degrees of freedom. Large-sample normal
approximations require justification. Running many tests and reporting only
the most striking interval changes the procedure's coverage claims.

\section{A control separates explanations}
Association is a property of a joint distribution. A causal claim asks what
would change under an intervention. Both temperature and ice-cream sales
can correlate with swimming incidents; a shared seasonal cause can generate
association without the proposed direct mechanism. A randomized comparison
balances causes in expectation under its design assumptions. Observational
causal inference needs additional assumptions about confounding and selection.

For simulations, matched initial conditions and recorded parameter changes
help isolate the mechanism being compared. Independent seed clusters matter
when outputs within one run share a trajectory. Report the difference between
models with uncertainty, rather than interpreting one model's attractive
picture as evidence of a specific cause. Chapter~\ref{learn:numerical}
adds implementation error to these statistical questions.

\section{Exercises with full solutions}
\paragraph{1. Change the population.}
Keep the diagnostic sensitivity and false-positive probability above, but
set prevalence to $0.10$. Find $P(D\mid+)$.

\paragraph{Solution.}
The positive probability is $0.99(0.10)+0.05(0.90)=0.144$.
The joint probability of the condition and a positive result is $0.099$,
so $P(D\mid+)=0.099/0.144=11/16=0.6875$.
The same test characteristics give a different posterior probability because
the population entering the calculation changed.

\paragraph{2. Compute an uncertainty scale.}
Independent Bernoulli trials have $p=1/2$. Find the standard deviation of
the sample mean for $n=100$ and $n=400$.

\paragraph{Solution.}
The formula gives $\sqrt{(1/2)(1/2)/100}=0.05$ and
$\sqrt{(1/2)(1/2)/400}=0.025$. Both numbers describe the sampling
spread under known $p=1/2$; estimating an unknown $p$ adds an inferential
step. Fourfold sample size halves the standard deviation.

\paragraph{3. Check whether two events are independent.}
For a fair die, let $A$ be an even result and $B$ a result at least four.

\paragraph{Solution.}
$P(A)=3/6=1/2$ and $P(B)=3/6=1/2$. Their intersection is
$\{4,6\}$, so $P(A\cap B)=1/3$, which differs from $P(A)P(B)=1/4$.
The events are dependent. Equivalently,
$P(A\mid B)=2/3\ne1/2=P(A)$.
